% =========================================================
% The Mediant Property
% =========================================================

\subsection{The Mediant Property}
\label{sec:mediant}

\subsubsection{Definitions and Theorems}

Once density is known abstractly, it is natural to ask for a \emph{constructive} way to produce a
new rational number between two old ones.  The mediant does exactly this.

\begin{definition}[Mediant]
\label{def:mediant}
Let
\[
\frac{a}{b},\frac{c}{d}\in\mathbb{Q},
\qquad b,d>0.
\]
The \emph{mediant} of these two fractions is
\[
\frac{a+c}{b+d}.
\]
\end{definition}

\begin{remark}
The mediant is formed by adding the numerators and denominators separately.  At first glance this
looks arbitrary, but it has a remarkable order-theoretic property.
\end{remark}

\begin{theorem}[Mediant inequality]
\label{thm:mediant-inequality}
If
\[
\frac{a}{b}<\frac{c}{d}
\qquad\text{with}\qquad b,d>0,
\]
then
\[
\frac{a}{b}<\frac{a+c}{b+d}<\frac{c}{d}.
\]
\end{theorem}

\begin{proof}
From $a/b<c/d$ we obtain $ad<bc$.
For the left-hand inequality, compute
\[
a(b+d)=ab+ad < ab+bc = b(a+c),
\]
so
\[
\frac{a}{b}<\frac{a+c}{b+d}.
\]
For the right-hand inequality,
\[
(a+c)d = ad+cd < bc+cd = c(b+d),
\]
which gives
\[
\frac{a+c}{b+d}<\frac{c}{d}.
\]
\end{proof}

\begin{corollary}[Density of $\mathbb{Q}$ within itself]
\label{cor:q-dense-in-q}
If $r,s\in\mathbb{Q}$ with $r<s$, then there exists $t\in\mathbb{Q}$ such that
\[
r<t<s.
\]
\end{corollary}

\begin{proof}
Write $r=a/b$ and $s=c/d$ with positive denominators.  Then the mediant
\[
t=\frac{a+c}{b+d}
\]
satisfies $r<t<s$.
\end{proof}

\begin{remark}[Local versus global structure]
\label{rem:mediant-vs-arch}
The mediant property and the Archimedean property describe different aspects of $\mathbb{Q}$.  The
mediant property is \emph{local}: it refines an interval by inserting a rational between two given
rationals.  The Archimedean property is \emph{global}: it says the natural numbers eventually exceed
any fixed number.
\end{remark}

\subsubsection{Farey and Stern--Brocot phenomena}

\begin{theorem}[Farey neighbor determinant condition]
\label{thm:farey-neighbor}
Suppose
\[
\frac{a}{b}<\frac{c}{d}
\qquad\text{with}\qquad
\gcd(a,b)=\gcd(c,d)=1,
\quad b,d>0.
\]
If
\[
bc-ad=1,
\]
then the two fractions are neighbors in a Farey sequence, and their mediant is already in reduced
form:
\[
\frac{a+c}{b+d}.
\]
\end{theorem}

\begin{remark}
The determinant condition $bc-ad=1$ is the arithmetic shadow of the Stern--Brocot tree.  Adjacent
fractions in that tree satisfy this condition, and the mediant inserts the unique next fraction between
them.
\end{remark}

\begin{theorem}[Stern--Brocot enumeration]
\label{thm:stern-brocot-enumeration}
Every positive rational number appears exactly once in reduced form in the Stern--Brocot tree.
\end{theorem}

\begin{remark}
This theorem is striking because it turns the set of positive rational numbers into a navigable object.
Repeated mediants do not merely produce some rationals; they eventually produce \emph{all} positive
rationals, and without duplication.
\end{remark}

\subsubsection{Consequences}

\begin{remark}
The mediant is one of the most useful constructive tools in elementary number theory and analysis.
It underlies Farey sequences, the Stern--Brocot tree, continued fractions, and many concrete
approximation procedures.
\end{remark}
